% F082｜03 三种证明：初中几何、有理化与不等式、高中导数
\clearpage
\begin{proofbox}
\textbf{\textcolor{CProof}{证明一：初中几何法（理解门槛最低）}}
\begin{center}\begin{tikzpicture}[font=\small,>=stealth,line join=round]
\draw[gray!40,rounded corners=3pt] (0,0) rectangle (14.2,8.4);\node at (7.1,7.800000000000001) {比较的是直路 PA 与绕路 PBA，先得到范围界};
\begin{scope}[shift={(3.0,2.2)},x=1.2cm,y=1.2cm]\draw[->,gray] (-1.6,0)--(2.6,0) node[right]{$x$};\draw[->,gray] (0,-.35)--(0,2.8) node[left]{$y$};\draw[orange!85!black,line width=1.6pt] (0,2)--(-1,0);\draw[violet,thick] (2,1)--(-1,0);\draw[teal!80!black,thick] (0,2)--(2,1);\fill[teal!80!black] (0,2) circle(2.2pt) node[above left=4pt]{$A(0,2)$};\fill[teal!80!black] (2,1) circle(2.2pt) node[above right=4pt]{$B(2,1)$};\fill[blue] (-1,0) circle(2.2pt) node[below=4pt]{$P(-1,0)$};\end{scope}\node[anchor=west,align=left] at (8.6,4.4) {直路：$PA$\\[12pt]绕行：$PB+BA$\\[12pt]$PA\le PB+BA$\\[12pt]$PA-PB\le AB$\\[12pt]交换 $A,B$ 再得另一侧界。};
\node[text width=12.8cm,align=center] at (7.1,.65) {橙色为直路；紫色与青色合起来是绕路。这里只得到界，还未证明能取到。};
\end{tikzpicture}\end{center}



\textbf{第 1 步：按勾股定理画点。} 取
\[
 A(m,a),\quad B(n,b),\quad P(x,0),
\]
则
\[
 PA=\sqrt{(x-m)^2+a^2},\quad
 PB=\sqrt{(x-n)^2+b^2},\quad
 f(x)=PA-PB.
\]

\textbf{第 2 步：先把差的范围圈住。} 从 $P$ 经 $B$ 到 $A$，总路程 $PB+BA$ 不可能短于 $PA$，所以 $PA-PB\leq AB$。交换 $A,B$ 可得 $PB-PA\leq AB$。因此
\[
 \boxed{-AB\leq f(x)\leq AB},\qquad
 AB=\sqrt{(m-n)^2+(a-b)^2}=D.
\]
\end{proofbox}
\clearpage
\begin{proofbox}
\textbf{第 3 步：看等号能否真正发生。} 要使差的绝对值等于 $AB$，$P,A,B$ 必须在\textbf{同一条直线上}，且\textbf{$P$ 在 $AB$ 的外侧}。
\begin{center}\begin{tikzpicture}[font=\small,>=stealth,line join=round]
\draw[gray!40,rounded corners=3pt] (0,0) rectangle (14.2,8.0);\node at (7.1,7.4) {高低换过来，延长线的交点与距离差的正负也改变};
\begin{scope}[shift={(2.0,3.2)},x=0.75cm,y=0.75cm]\draw[->,gray] (-0.6,0)--(4.6,0) node[right]{$x$};\draw[->,gray] (0,-.35)--(0,2.8) node[left]{$y$};\draw[orange!85!black,line width=1.6pt] (0,2)--(4,0);\draw[violet,thick] (2,1)--(4,0);\draw[teal!80!black,thick] (0,2)--(2,1);\fill[teal!80!black] (0,2) circle(2.2pt) node[left=6pt]{$A$};\fill[teal!80!black] (2,1) circle(2.2pt) node[above right=4pt]{$B$};\fill[blue] (4,0) circle(2.2pt) node[below=4pt]{$P$};\end{scope}\node at (3.2,6.3) {$a>b$：正端点};\node at (3.2,1.65) {$f(4)=\sqrt5$};\begin{scope}[shift={(10.0,3.2)},x=0.75cm,y=0.75cm]\draw[->,gray] (-2.6,0)--(2.6,0) node[right]{$x$};\draw[->,gray] (0,-.35)--(0,2.8) node[left]{$y$};\draw[orange!85!black,line width=1.6pt] (0,1)--(-2,0);\draw[violet,thick] (2,2)--(-2,0);\draw[teal!80!black,thick] (0,1)--(2,2);\fill[teal!80!black] (0,1) circle(2.2pt) node[left=6pt]{$A$};\fill[teal!80!black] (2,2) circle(2.2pt) node[above right=4pt]{$B$};\fill[blue] (-2,0) circle(2.2pt) node[below=4pt]{$P$};\end{scope}\node at (11.2,6.3) {$a<b$：负端点};\node at (11.2,1.65) {$f(-2)=-\sqrt5$};
\node[text width=12.8cm,align=center] at (7.1,.65) {同尺度对照：$a\ne b$ 时 $x_0=(bm-an)/(b-a)$；交点都在线段外侧。};
\end{tikzpicture}\end{center}



若 $a\ne b$，两点高度不一样，$AB$ 的延长线一定会与 $x$ 轴相交。交点 $P_0$ 满足
\[
 \frac{x_0-m}{n-m}=\frac{-a}{b-a}\quad(n\ne m).
\]
解得
\[
 \boxed{x_0=\frac{bm-an}{b-a}}.
\]
当 $m=n$ 时，直接看竖直线 $x=m$，仍得同一公式。若 $a>b$，$P_0$ 在 $B$ 的外侧，所以 $PA-PB=AB=D$；若 $a<b$，$P_0$ 在 $A$ 的外侧，所以 $PA-PB=-AB=-D$。

\end{proofbox}
\clearpage
\begin{proofbox}
\textbf{\textcolor{CProof}{证明一（续）：高度相同时为什么等号不成立？}}

若 $a=b$、$m\ne n$，$AB$ 是一条位于 $x$ 轴上方的水平线。无论 $P$ 在 $x$ 轴上怎么移动，都不能落到这条线上，因此
\[
 -|m-n|<PA-PB<|m-n|.
\]
\begin{center}\begin{tikzpicture}[font=\small,>=stealth,line join=round]
\draw[gray!40,rounded corners=3pt] (0,0) rectangle (14.2,8.0);\node at (7.1,7.4) {两定点一样高：延长线与 x 轴永远不相交};
\begin{scope}[shift={(2.7,2.3)},x=1cm,y=1cm]\draw[->,gray] (-1.6,0)--(3.6,0) node[right]{$x$};\draw[->,gray] (0,-.35)--(0,2.8) node[left]{$y$};\draw[orange!85!black,line width=1.6pt] (0,2)--(-1,0);\draw[violet,thick] (3,2)--(-1,0);\draw[teal!80!black,thick] (0,2)--(3,2);\fill[teal!80!black] (0,2) circle(2.2pt) node[above left=4pt]{$A(0,2)$};\fill[teal!80!black] (3,2) circle(2.2pt) node[above right=4pt]{$B(3,2)$};\fill[blue] (-1,0) circle(2.2pt) node[below=4pt]{$P(-1,0)$};\end{scope}\draw[teal!80!black,dashed] (1.2,4.3)--(8,4.3);\node[anchor=west,align=left] at (8.6,4.4) {$A(0,2),\ B(3,2)$\\[12pt]$AB$ 所在直线为 $y=2$\\[12pt]与动点所在 $y=0$ 平行\\[12pt]不存在取等所需的交点。};
\node[text width=12.8cm,align=center] at (7.1,.65) {只要两定点不同且都在轴上方，距离差的两端都不能取到。};
\end{tikzpicture}\end{center}


\textbf{第 4 步：看走得特别远时会怎样。} 当 $x$ 很大时，根号 $\sqrt{(x-m)^2+a^2}$ 与 $|x-m|$ 的差会越来越小。理由只需有理化：
\[
 \sqrt{(x-m)^2+a^2}-|x-m|
 =\frac{a^2}{\sqrt{(x-m)^2+a^2}+|x-m|}\longrightarrow0.
\]
第二个根号同理。所以
\[
 x\to+\infty:\quad f(x)\to n-m,\qquad
 x\to-\infty:\quad f(x)\to m-n.
\]
由两端可以无限接近的值、刚才证明的取等位置，再配合函数值随 $x$ 连续变化，可确定所有中间数都能取到。

不过，\textbf{还有一步不能省：}必须证明距离差不会在某处穿过“另一端”的 $\pm d$。这个严格判断在证明二中用初中平方与有理化完成。最终结果是：
\[
 a>b:\ (-d,D],\quad a<b:\ [-D,d),\quad
 a=b,\ m\ne n:\ (-d,d).
\]
若 $a=b,m=n$，则原式一直为 $0$。
\end{proofbox}
\clearpage
\begin{proofbox}
\textbf{\textcolor{CProof}{证明二：有理化与平方不等式法（用图认清辅助量，再展开计算）}}
\begin{center}\begin{tikzpicture}[font=\small,>=stealth,line join=round]
\draw[gray!40,rounded corners=3pt] (0,0) rectangle (14.2,8.8);\node at (7.1,8.200000000000001) {平方推导中的 U、V、D，就是同一三角形的三条边};
\begin{scope}[shift={(4.4,2.4)},x=1.3cm,y=1.3cm]\draw[->,gray] (-1.7,0)--(1.7,0) node[right]{$t\text{方向}$};\draw[->,gray] (0,-.3)--(0,2.8);\draw[orange!85!black,thick,->] (0,0)--(1,2);\draw[violet,thick,->] (0,0)--(-1,1);\draw[teal!80!black,thick,->] (-1,1)--(1,2);\fill (0,0) circle(2pt) node[below right]{$O$};\fill (1,2) circle(2pt) node[above right=3pt]{$C(t,a)$};\fill (-1,1) circle(2pt) node[above left=3pt]{$E(s,b)$};\end{scope}\node[anchor=west,align=left] at (8.8,4.2) {$OC=U=\sqrt{t^2+a^2}$\\[11pt]$OE=V=\sqrt{s^2+b^2}$\\[11pt]$CE=D$\\[12pt]$\lvert U-V\rvert\le D$\\[11pt]本例 $(bt-as)^2=9$};
\node[text width=12.8cm,align=center] at (7.1,.65) {图取 $t=1,s=-1,a=2,b=1$；$CE^2=(t-s)^2+(a-b)^2$。};
\end{tikzpicture}\end{center}



\textbf{第一部分：先证明差的绝对值不超过 $D$。} 设
\[
 t=x-m,\quad s=x-n,\quad
 U=\sqrt{t^2+a^2},\quad V=\sqrt{s^2+b^2}.
\]
由于任何平方都不小于 $0$，直接展开可得
\[
 \begin{aligned}
 (UV)^2-(ts+ab)^2
 &=(t^2+a^2)(s^2+b^2)-(ts+ab)^2\\
 &=b^2t^2+a^2s^2-2abts\\
 &=(bt-as)^2\geq0.
 \end{aligned}
\]
因此 $UV\geq|ts+ab|\geq ts+ab$。所以
\[
 \begin{aligned}
 (U-V)^2
 &=t^2+a^2+s^2+b^2-2UV\\
 &\leq t^2+s^2-2ts+a^2+b^2-2ab\\
 &=(t-s)^2+(a-b)^2\\
 &=(n-m)^2+(a-b)^2=D^2.
 \end{aligned}
\]
故 $\boxed{-D\leq f(x)\leq D}$。

\end{proofbox}
\clearpage
\begin{proofbox}
\textbf{第二部分：哪个端点能取到？} 上面用到的平方等于 $0$，需要
\begin{center}\begin{tikzpicture}[font=\small,>=stealth,line join=round]
\draw[gray!40,rounded corners=3pt] (0,0) rectangle (14.2,8.0);\node at (7.1,7.4) {高低换过来，延长线的交点与距离差的正负也改变};
\begin{scope}[shift={(2.0,3.2)},x=0.75cm,y=0.75cm]\draw[->,gray] (-0.6,0)--(4.6,0) node[right]{$x$};\draw[->,gray] (0,-.35)--(0,2.8) node[left]{$y$};\draw[orange!85!black,line width=1.6pt] (0,2)--(4,0);\draw[violet,thick] (2,1)--(4,0);\draw[teal!80!black,thick] (0,2)--(2,1);\fill[teal!80!black] (0,2) circle(2.2pt) node[left=6pt]{$A$};\fill[teal!80!black] (2,1) circle(2.2pt) node[above right=4pt]{$B$};\fill[blue] (4,0) circle(2.2pt) node[below=4pt]{$P$};\end{scope}\node at (3.2,6.3) {$a>b$：正端点};\node at (3.2,1.65) {$f(4)=\sqrt5$};\begin{scope}[shift={(10.0,3.2)},x=0.75cm,y=0.75cm]\draw[->,gray] (-2.6,0)--(2.6,0) node[right]{$x$};\draw[->,gray] (0,-.35)--(0,2.8) node[left]{$y$};\draw[orange!85!black,line width=1.6pt] (0,1)--(-2,0);\draw[violet,thick] (2,2)--(-2,0);\draw[teal!80!black,thick] (0,1)--(2,2);\fill[teal!80!black] (0,1) circle(2.2pt) node[left=6pt]{$A$};\fill[teal!80!black] (2,2) circle(2.2pt) node[above right=4pt]{$B$};\fill[blue] (-2,0) circle(2.2pt) node[below=4pt]{$P$};\end{scope}\node at (11.2,6.3) {$a<b$：负端点};\node at (11.2,1.65) {$f(-2)=-\sqrt5$};
\node[text width=12.8cm,align=center] at (7.1,.65) {同尺度对照：$a\ne b$ 时 $x_0=(bm-an)/(b-a)$；交点都在线段外侧。};
\end{tikzpicture}\end{center}


\[
 bt=as\quad\Longleftrightarrow\quad
 b(x-m)=a(x-n).
\]
若 $a\ne b$，解得
\[
 \boxed{x_0=\frac{bm-an}{b-a}}.
\]
此时 $t/a=s/b$，设其共同数为 $r$，则
\[
 U=a\sqrt{1+r^2},\quad V=b\sqrt{1+r^2},\quad
 U-V=(a-b)\sqrt{1+r^2}.
\]
又 $D=|a-b|\sqrt{1+r^2}$，因此 $a>b$ 时恰有 $f(x_0)=D$，$a<b$ 时恰有 $f(x_0)=-D$。

若 $a=b$ 且 $m\ne n$，等号条件 $t=s$ 不可能成立；若 $a=b,m=n$，原式恒为 $0$。
\end{proofbox}
\clearpage
\begin{proofbox}
\textbf{\textcolor{CProof}{证明二（续）：用“有理化”判断另一端一定取不到}}
\begin{center}\begin{tikzpicture}[font=\small,>=stealth,line join=round]
\draw[gray!40,rounded corners=3pt] (0,0) rectangle (14.2,8.7);\node at (7.1,8.1) {降低第一个定点：多出的正距离使下边界取不到};
\begin{scope}[shift={(3,2.3)},x=1.15cm,y=1.15cm]\draw[->,gray] (-1.6,0)--(2.6,0) node[right]{$x$};\draw[->,gray] (0,-.35)--(0,2.8) node[left]{$y$};\draw[orange!85!black,line width=1.6pt] (0,2)--(-1,0);\draw[violet,thick] (2,1)--(-1,0);\draw[teal!80!black,thick] (0,2)--(2,1);\fill[teal!80!black] (0,2) circle(2.2pt) node[above left=4pt]{$A(0,2)$};\fill[teal!80!black] (2,1) circle(2.2pt) node[above right=4pt]{$B(2,1)$};\fill[blue] (-1,0) circle(2.2pt) node[below=4pt]{$P(-1,0)$};\end{scope}\begin{scope}[shift={(3,2.3)},x=1.15cm,y=1.15cm]\draw[gray,thick,dashed] (-1,0)--(0,1);\draw[red!70!black,->] (0,2)--(0,1);\fill[gray] (0,1) circle(2pt) node[left=4pt]{$C$};\end{scope}\node[anchor=west,align=left] at (8.7,4.3) {$W=PC$\\[12pt]$U-W>0$\\[12pt]$f=(U-W)+(W-V)$\\[12pt]$W-V>-d$\\[12pt]所以 $f>-d$。};
\node[text width=12.8cm,align=center] at (7.1,.65) {灰虚线是降低高度后的距离 $W$；红箭头改变固定高度，不改变动点所在的轴。};
\end{tikzpicture}\end{center}



\textbf{情形一：$a>b$。} 把第一个根号的高度暂时改成 $b$：
\[
 W=\sqrt{(x-m)^2+b^2}.
\]
由于 $a>b$，第一个根号 $U$ 严格大于 $W$，并且
\[
 \underbrace{U-W}_{\text{严格大于 }0}
 =\frac{a^2-b^2}{U+W}>0.
\]
原式拆成
\[
 f(x)=\underbrace{(U-W)}_{>0}+\underbrace{(W-V)}_{\text{两个同高度点的距离差}}.
\]
第二项直接有理化：
\[
 W-V=\frac{(x-m)^2-(x-n)^2}{W+V}
 =\frac{(n-m)(2x-m-n)}{W+V}.
\]
若 $m\ne n$，因为 $b>0$，有 $W>|x-m|$、$V>|x-n|$，从而
\[
 |W-V|
 =|n-m|\,\frac{|(x-m)+(x-n)|}{W+V}
 <|n-m|=d.
\]
若 $m=n$，直接有 $W-V=0$。所以无论 $m,n$ 是否相等，都有
\[
 \boxed{f(x)>-d}.
\]
当 $x$ 往足够远的一端移动时，$U-W$ 越来越接近 $0$，而 $W-V$ 越来越接近 $-d$，因此 $-d$ 是一个\textbf{只能靠近、不能取得}的边界。

综合 $f(x_0)=D$、$f(x)>-d$、函数的连续变化，得到
\[
 \boxed{a>b:\qquad f(x)\in(-d,D]}.
\]
\end{proofbox}
\clearpage
\begin{proofbox}
\textbf{\textcolor{CProof}{证明二（续）：$a<b$ 和 $a=b$}}
\begin{center}\begin{tikzpicture}[font=\small,>=stealth,line join=round]
\draw[gray!40,rounded corners=3pt] (0,0) rectangle (14.2,8.7);\node at (7.1,8.1) {降低第二个定点：减掉正距离使上边界取不到};
\begin{scope}[shift={(1.9,2.3)},x=1.15cm,y=1.15cm]\draw[->,gray] (-0.6,0)--(3.6,0) node[right]{$x$};\draw[->,gray] (0,-.35)--(0,2.8) node[left]{$y$};\draw[orange!85!black,line width=1.6pt] (0,1)--(3,0);\draw[violet,thick] (2,2)--(3,0);\draw[teal!80!black,thick] (0,1)--(2,2);\fill[teal!80!black] (0,1) circle(2.2pt) node[above left=4pt]{$A(0,1)$};\fill[teal!80!black] (2,2) circle(2.2pt) node[above right=4pt]{$B(2,2)$};\fill[blue] (3,0) circle(2.2pt) node[below=4pt]{$P(3,0)$};\end{scope}\begin{scope}[shift={(1.9,2.3)},x=1.15cm,y=1.15cm]\draw[gray,thick,dashed] (3,0)--(2,1);\draw[red!70!black,->] (2,2)--(2,1);\fill[gray] (2,1) circle(2pt) node[left=4pt]{$C$};\end{scope}\node[anchor=west,align=left] at (8.7,4.3) {$W=PC$\\[12pt]$V-W>0$\\[12pt]$f=(U-W)-(V-W)$\\[12pt]$U-W\le d$\\[12pt]所以 $f<d$。};
\node[text width=12.8cm,align=center] at (7.1,.65) {灰虚线是降低高度后的距离 $W$；红箭头改变固定高度，不改变动点所在的轴。};
\end{tikzpicture}\end{center}



\textbf{情形二：$a<b$。} 这回把第二个根号的高度暂时\textbf{降低}到 $a$。令
\[
 W=\sqrt{(x-n)^2+a^2}.
\]
由于 $V>W$，有理化给出
\[
 V-W=\frac{b^2-a^2}{V+W}>0.
\]
于是
\[
 f(x)=(U-W)-(V-W).
\]
其中 $U,W$ 对应的两个定点高度同为 $a>0$，同样由有理化知道
\[
 U-W\leq d \quad\text{（$m=n$ 时可能等于 $d=0$）}.
\]
再减掉严格为正的 $V-W$，便得到
\[
 \boxed{f(x)<d}.
\]
前面已证明 $f(x_0)=-D$，而往远处移动时可以无限接近 $d$，故
\[
 \boxed{a<b:\qquad f(x)\in[-D,d)}.
\]

\end{proofbox}
\clearpage
\begin{proofbox}
\begin{center}\begin{tikzpicture}[font=\small,>=stealth,line join=round]
\draw[gray!40,rounded corners=3pt] (0,0) rectangle (14.2,9.2);\node at (7.1,8.6) {高度相同：一直增大，两端都只能靠近};
\node[blue] at (7.1,7.65) {$\displaystyle h(x)=\sqrt{x^2+4}-\sqrt{(x-3)^2+4}$};\begin{scope}[shift={(4.1,4.0)},x=0.3130434782608696cm,y=0.75cm]\draw[->,gray] (-10,0)--(13.5,0) node[right]{$x$};\draw[->,gray] (0,-3.6)--(0,3.8000000000000003) node[left]{$y$};\draw[gray,dashed] (-10,-3)--(13,-3);\node[left=3pt,fill=white,inner sep=1pt] at (0,-3) {$-3$};\draw[gray,dashed] (-10,3)--(13,3);\node[left=3pt,fill=white,inner sep=1pt] at (0,3) {$3$};\draw[blue,very thick,domain=-10:13,samples=240] plot(\x,{sqrt(\x*\x+4)-sqrt((\x-3)^2+4)});\end{scope}\node[anchor=west,align=left] at (9,4.5) {$a=b=2,\ m=0,n=3$\\[12pt]对每个实数 $x$\\[10pt]都有 $-3<h(x)<3$\\[12pt]$h(3/2)=0$\\[12pt]完整值域：$(-3,3)$};
\node[text width=12.8cm,align=center] at (7.1,.65) {两条虚线均不属于函数图像；同高度且不同中心时没有最大、最小值。};
\end{tikzpicture}\end{center}

\textbf{情形三：$a=b$。} 此时两个根号的固定高度一样：
\[
 f(x)=\frac{(n-m)(2x-m-n)}{
 \sqrt{(x-m)^2+a^2}+\sqrt{(x-n)^2+a^2}}.
\]
若 $m\ne n$，依照上面的分母比较，严格有 $|f(x)|<d$；往左右无穷远分别靠近 $-d$ 与 $d$，所以
\[
 \boxed{a=b,\ m\ne n:\quad f(x)\in(-d,d)}.
\]
若 $m=n$，分子恒为 $0$，于是 $f(x)\equiv0$，值域为 $\{0\}$。

\textbf{本证明最有用的想法：}一个高度大一点，就把它\textbf{降到另一个的高度}；多出来的那一截距离严格为正，而“同高度的距离差”永远小于水平距离。这样不用猜，便能证明哪个端点取不到。
\end{proofbox}
\clearpage
\begin{proofbox}
\textbf{\textcolor{CProof}{证明三：高中导数法（常规方法，可处理限制区间）}}
\begin{center}\begin{tikzpicture}[font=\small,>=stealth,line join=round]
\draw[gray!40,rounded corners=3pt] (0,0) rectangle (14.2,8.6);\node at (7.1,8.0) {两个复杂分式的大小，转成比较输入的大小};
\begin{scope}[shift={(4.0,4.2)},x=1cm,y=2.0cm]\draw[->,gray] (-3.1,0)--(3.2,0) node[right]{$t$};\draw[->,gray] (0,-1.2)--(0,1.25) node[left]{$q$};\draw[gray,dashed] (-3,1)--(3,1) (-3,-1)--(3,-1);\draw[blue,thick,domain=-3:3,samples=130] plot(\x,{\x/sqrt(1+\x*\x)});\fill[orange!85!black] (-1,{-1/sqrt(2)}) circle(2pt);\fill[orange!85!black] (1,{1/sqrt(2)}) circle(2pt);\end{scope}\node[anchor=west,align=left] at (8.5,4.3) {$q(t)=\dfrac{t}{\sqrt{1+t^2}}$\\[14pt]$q$ 严格递增\\[12pt]比较 $q(r)-q(s)$ 的符号\\[10pt]只需比较 $r-s$ 的符号\\[12pt]这里 $r=(x-m)/a$\\[8pt]$s=(x-n)/b$。};
\node[text width=12.8cm,align=center] at (7.1,.65) {橙点示意 $-1<1$ 时 $q(-1)<q(1)$；全区间严格递增由 $q\prime(t)>0$ 证明。};
\end{tikzpicture}\end{center}



记
\[
 U(x)=\sqrt{(x-m)^2+a^2},\quad
 V(x)=\sqrt{(x-n)^2+b^2},\quad f(x)=U(x)-V(x).
\]
直接求导：
\[
 f'(x)=\frac{x-m}{\sqrt{(x-m)^2+a^2}}
       -\frac{x-n}{\sqrt{(x-n)^2+b^2}}.
\]
两个分式看似复杂，关键是：对任意实数 $t$，
\[
 q(t)=\frac{t}{\sqrt{1+t^2}},\qquad
 q'(t)=\frac{1}{(1+t^2)^{3/2}}>0.
\]
所以 $q(t)$ 随 $t$ 增大而增大。原来的两个分式恰好是
\[
 f'(x)=q\!\left(\frac{x-m}{a}\right)
       -q\!\left(\frac{x-n}{b}\right).
\]
因此比较 $f'(x)$ 的正负，只需比较括号里的两个数：
\[
 \begin{aligned}
 f'(x)>0
 &\Longleftrightarrow \frac{x-m}{a}>\frac{x-n}{b}\\
 &\Longleftrightarrow b(x-m)>a(x-n)\\
 &\Longleftrightarrow (b-a)x>bm-an.
 \end{aligned}
\]
把“大于”换成“等于”或“小于”，也完全成立。

\end{proofbox}
\clearpage
\begin{proofbox}
\textbf{当 $a>b$ 时：} $b-a<0$，以 $x_0=(bm-an)/(b-a)$ 为界，
\begin{center}\begin{tikzpicture}[font=\small,>=stealth,line join=round]
\draw[gray!40,rounded corners=3pt] (0,0) rectangle (14.2,10.2);\node at (7.1,9.6) {第一项高度更大：先增后减，只有最大值能取到};
\node[blue] at (7.1,8.65) {$\displaystyle f(x)=\sqrt{x^2+4}-\sqrt{(x-2)^2+1}$};\begin{scope}[shift={(3.1,4.6)},x=0.36cm,y=1cm]\draw[->,gray] (-6,0)--(14.5,0) node[right]{$x$};\draw[->,gray] (0,-2.8)--(0,3.4000000000000004) node[left]{$y$};\draw[gray,dashed] (-6,-2)--(14,-2);\node[left=3pt,fill=white,inner sep=1pt] at (0,-2) {$-2$};\draw[gray,dashed] (-6,2)--(14,2);\node[left=3pt,fill=white,inner sep=1pt] at (0,2) {$2$};\draw[blue,very thick,domain=-6:14,samples=240] plot(\x,{sqrt(\x*\x+4)-sqrt((\x-2)^2+1)});\draw[gray,dotted] (4,0)--(4,2.23606797749979);\fill[blue] (4,2.23606797749979) circle(2.3pt);\node[below=4pt] at (4,0){$4$};\node[blue,above=9pt] at (4,2.23606797749979) {$(4,\sqrt5)$};\end{scope}\node[anchor=west,align=left] at (9,4.5) {$a=2>b=1$\\[12pt]左端：趋近 $-2$\\[10pt]右端：趋近 $2$\\[12pt]最大值：$\sqrt5$\\[12pt]完整值域：$(-2,\sqrt5]$};
\node[text width=12.8cm,align=center] at (7.1,.65) {虚线是无限远处靠近的高度；不是有限定义域的端点。};
\end{tikzpicture}\end{center}


\[
 f'(x)\begin{cases}>0,&x<x_0,\\=0,&x=x_0,\\<0,&x>x_0.\end{cases}
\]
故函数\textbf{先增大后减小}，在 $x_0$ 处取到最大值 $D$。
\end{proofbox}
\clearpage
\begin{proofbox}
\textbf{\textcolor{CProof}{证明三（续）：另外两种情况与完整值域}}
\begin{center}\begin{tikzpicture}[font=\small,>=stealth,line join=round]
\draw[gray!40,rounded corners=3pt] (0,0) rectangle (14.2,10.2);\node at (7.1,9.6) {第一项高度更小：先减后增，只有最小值能取到};
\node[blue] at (7.1,8.65) {$\displaystyle g(x)=\sqrt{x^2+1}-\sqrt{(x-2)^2+4}$};\begin{scope}[shift={(3.1,4.6)},x=0.36cm,y=1cm]\draw[->,gray] (-6,0)--(14.5,0) node[right]{$x$};\draw[->,gray] (0,-2.8)--(0,3.4000000000000004) node[left]{$y$};\draw[gray,dashed] (-6,-2)--(14,-2);\node[left=3pt,fill=white,inner sep=1pt] at (0,-2) {$-2$};\draw[gray,dashed] (-6,2)--(14,2);\node[left=3pt,fill=white,inner sep=1pt] at (0,2) {$2$};\draw[blue,very thick,domain=-6:14,samples=240] plot(\x,{sqrt(\x*\x+1)-sqrt((\x-2)^2+4)});\draw[gray,dotted] (-2,0)--(-2,-2.23606797749979);\fill[blue] (-2,-2.23606797749979) circle(2.3pt);\node[below=4pt] at (-2,0){$-2$};\node[blue,below=9pt] at (-2,-2.23606797749979) {$(-2,-\sqrt5)$};\end{scope}\node[anchor=west,align=left] at (9,4.5) {$a=1<b=2$\\[12pt]左端：趋近 $-2$\\[10pt]右端：趋近 $2$\\[12pt]最小值：$-\sqrt5$\\[12pt]完整值域：$[-\sqrt5,2)$};
\node[text width=12.8cm,align=center] at (7.1,.65) {橙紫两项顺序固定；这里的最低点为负，不能把距离差改成绝对值。};
\end{tikzpicture}\end{center}



\textbf{当 $a<b$ 时：} $b-a>0$，导数在 $x_0$ 左边为负、右边为正，函数\textbf{先减小后增大}，在 $x_0$ 处取到最小值 $-D$。

\end{proofbox}
\clearpage
\begin{proofbox}
\textbf{当 $a=b$ 时：} 比较式变成
\begin{center}\begin{tikzpicture}[font=\small,>=stealth,line join=round]
\draw[gray!40,rounded corners=3pt] (0,0) rectangle (14.2,9.2);\node at (7.1,8.6) {高度相同：一直增大，两端都只能靠近};
\node[blue] at (7.1,7.65) {$\displaystyle h(x)=\sqrt{x^2+4}-\sqrt{(x-3)^2+4}$};\begin{scope}[shift={(4.1,4.0)},x=0.3130434782608696cm,y=0.75cm]\draw[->,gray] (-10,0)--(13.5,0) node[right]{$x$};\draw[->,gray] (0,-3.6)--(0,3.8000000000000003) node[left]{$y$};\draw[gray,dashed] (-10,-3)--(13,-3);\node[left=3pt,fill=white,inner sep=1pt] at (0,-3) {$-3$};\draw[gray,dashed] (-10,3)--(13,3);\node[left=3pt,fill=white,inner sep=1pt] at (0,3) {$3$};\draw[blue,very thick,domain=-10:13,samples=240] plot(\x,{sqrt(\x*\x+4)-sqrt((\x-3)^2+4)});\end{scope}\node[anchor=west,align=left] at (9,4.5) {$a=b=2,\ m=0,n=3$\\[12pt]对每个实数 $x$\\[10pt]都有 $-3<h(x)<3$\\[12pt]$h(3/2)=0$\\[12pt]完整值域：$(-3,3)$};
\node[text width=12.8cm,align=center] at (7.1,.65) {两条虚线均不属于函数图像；同高度且不同中心时没有最大、最小值。};
\end{tikzpicture}\end{center}


\[
 b(x-m)-a(x-n)=a(n-m).
\]
如果 $m\ne n$，这个数始终同号，函数始终\textbf{只朝一个方向变化}，没有最大值也没有最小值；如果 $m=n$，导数恒为 $0$，函数恒等于 $0$。

\textbf{还需补上值域两端。} 当 $x$ 越来越大或越来越小时，
\[
 \lim_{x\to+\infty}f(x)=n-m,\qquad
 \lim_{x\to-\infty}f(x)=m-n.
\]
证明见“证明一”的有理化式。两端极限为 $d$ 与 $-d$（$m=n$ 时都为 $0$）。

导数已说明：$a>b$ 时只有一个最高点，其余两侧沿曲线向左右远处走，最低只会逼近 $-d$；$a<b$ 时只有一个最低点，最高只会逼近 $d$；$a=b$ 且 $m\ne n$ 时始终朝一个方向变化，两个端点都只能靠近。

因此完整值域是
\[
 \boxed{\begin{array}{c|c}
 a>b&(-d,D]\\
 a<b&[-D,d)\\
 a=b,\ m\ne n&(-d,d)\\
 a=b,\ m=n&\{0\}
 \end{array}}
\]
其中 $D=\sqrt{(m-n)^2+(a-b)^2}$，$d=|m-n|$。

\textbf{导数法适合什么时候？} 受限区间先检查 $x_0$，再比较允许的端点值。\textbf{端点未取到时，不能把对应函数值当作最值。}
\end{proofbox}
